% Source: https://artofproblemsolving.com/community/c3h1295244_polygon_angles
% Stable identity: aops-topic-1295244-q01
Let the smaller angle be \(x\); the larger is \(3x\). The figure has six small angles (at the arm ends) and three large (reflex) angles where the arms meet.

\begin{asy}
usepackage("amsmath");
size(4.5cm);
defaultpen(fontsize(10pt));
pair[] V; pair p=(0,0); real hd=90; real[] L={2.2,1.2,2.2}; real[] T={-72,96,96};
// walk: at each vertex turn, then walk the next edge; pattern reflex, small, small
real ang=hd; V.push(p);
for(int k=0;k<9;++k){ ang+=T[k%3]; p=p+L[k%3]*dir(ang); V.push(p);}
path g=V[0]; for(int k=1;k<9;++k) g=g--V[k]; g=g--cycle;
draw(g,linewidth(0.9));
for(int k=0;k<9;++k){ if(k%3==0) dot(V[k],red+linewidth(3)); else dot(V[k],blue+linewidth(3)); }
label("$3x$",V[0]+(-0.55,0.35),red); label("$x$",V[1],NE,blue); label("$x$",V[2],NW,blue);
label("\mbox{}",(min(currentpicture,user=true)+max(currentpicture,user=true))/2,black); // keeps the text after the figure black when LaTeX typesets the labels inline
\end{asy}

The interior angles of any nine-sided polygon add to \((9-2)\cdot180^\circ=1260^\circ\). Therefore
\[
6x+3(3x)=15x=1260^\circ,\qquad x=84^\circ,
\]
and each larger angle is \(3x=\boxed{252^\circ}\) (a reflex angle, as expected where the arms meet).
